%%

%

%%%%

%%%   Самарский государственный технический университет

%%%   Кафедра прикладной математики и информатики

%%%

%%%   Преамбула для статьи в журнал "Вестник Самарского государственного

%%%   технического университета. Серия: «Физико-математические науки»"

%%%   Ориентирована под MikTEx 2.4 for Win32

%%%

%%%   Хотя сам журнал собирается под GNU/Linux на дистрибутиве TeXLive



\documentclass[11pt,a4paper,twoside]{article}



\usepackage[cp1251]{inputenc}

\usepackage[T2A]{fontenc}

\usepackage[english,russian]{babel}

\usepackage{samgtubib}

\usepackage[dvips]{graphicx}

\usepackage[multiple]{footmisc}



\usepackage{samgtu}

\renewcommand{\VestnikYear}{2009} % Год издания Вестника

\renewcommand{\VestnikNumber}{2\,(19)} % Номер Вестника

\renewcommand{\VestnikMonth}{Ноябрь} % Месяц выхода Вестника

\renewcommand{\VestnikData}{01.11.09} % Дата подписания Вестника в печать

\renewcommand{\VestnikUPL}{26,64} % Усл. печ. л.

\renewcommand{\VestnikUIL}{25,73} % Уч.-изд. л.

\renewcommand{\VestnikRegNumber}{479/09} % Регистрационный номер

\renewcommand{\VestnikOrder}{994} % Номер заказа







%

%%%%%

%%%   Самарский государственный технический университет

%%%   Кафедра прикладной математики и информатики

%%%

%%%   Преамбула для статьи в журнал "Вестник Самарского государственного

%%%   технического университета. Серия: «Физико-математические науки»"

%%%   Ориентирована под MikTEx 2.4 for Win32

%%%

%%%   Хотя сам журнал собирается под GNU/Linux на дистрибутиве TeXLive



\documentclass[11pt,a4paper,twoside]{article}



\usepackage[cp1251]{inputenc}

\usepackage[T2A]{fontenc}

\usepackage[english,russian]{babel}

\usepackage{samgtubib}

\usepackage[dvips]{graphicx}

\usepackage[multiple]{footmisc}



\usepackage{samgtu}

\renewcommand{\VestnikYear}{2009} % Год издания Вестника

\renewcommand{\VestnikNumber}{2\,(19)} % Номер Вестника

\renewcommand{\VestnikMonth}{Ноябрь} % Месяц выхода Вестника

\renewcommand{\VestnikData}{01.11.09} % Дата подписания Вестника в печать

\renewcommand{\VestnikUPL}{26,64} % Усл. печ. л.

\renewcommand{\VestnikUIL}{25,73} % Уч.-изд. л.

\renewcommand{\VestnikRegNumber}{479/09} % Регистрационный номер

\renewcommand{\VestnikOrder}{994} % Номер заказа







%

%\usepackage{qrcode}

%

%%

%%\newcommand{\totop}[1]{\raisebox{6pt}[1pt][2pt]{#1}} 

%%\newcommand{\tottop}[1]{\raisebox{12pt}[1pt][2pt]{#1}} 

%%\newcommand{\totttop}[1]{\raisebox{18pt}[1pt][2pt]{#1}} 

%%\newcommand{\tottttop}[1]{\raisebox{24pt}[1pt][2pt]{#1}} 

%%\newcommand{\totttttop}[1]{\raisebox{30pt}[1pt][2pt]{#1}} 

%\newcommand{\tobot}[1]{\raisebox{-6pt}[1pt][2pt]{#1}} 

%%\newcommand{\tobbot}[1]{\raisebox{-12pt}[1pt][2pt]{#1}} 

%%\newcommand{\tobbbot}[1]{\raisebox{-18pt}[1pt][2pt]{#1}}

%

%\graphicspath{{./00PIC/}}

%

%%\samgtupubl % для генерации pdf, отдаваемого в типографию СамГТУ

%\pdfcrop % обрезка pdf для размещения на math-net.ru

%\draft % На корректуре

%\filllastpage % Последняя страница не заполнена не полностью

%%\briefcomm % Статья является кратким сообщением



\vsgtu{1493}



\FirstPage{544}

\LastPage{551}

%%

%%

%\begin{document}





%\renewcommand{\baselinestretch}{0.9}

%

%\hfill \qrcode[hyperlink,height=.8in]{http://dx.doi.org/10.14498/vsgtu1493}

%\vspace{-.8in}



% Номер УДК

\udc{519.6}%Вычислительная математика, численный анализ

\msc{97N40}% Коды MSC см. http://www.ams.org/msc/



\rutitle{Блочный регуляризованный метод Качмажа}{Блочный регуляризованный метод Качмажа}



\ruauthor{Е.~Ю.~Богданова}{Б\,о\,г\,д\,а\,н\,о\,в\,а~Е.~Ю.}% Автор(ы) для оглавления и колонтитула через \,



\ruaffil{\rusamgtu.}

\enaffil{\ensamgtu.}



%\email{E-mails}{fwinter@yandex.ru} %E-mail's авторов



\ruauthordetails{1}{\fullauthor{\rupid{119144}{Екатерина Юрьевна Богданова}}\regalia{\mem{fwinter@yandex.ru}} \position{аспирант} \dept{каф. высшей математики и прикладной информатики}}

\enauthordetails{1}{\fullauthor{\enpid{119144}{Ekaterina Yu. Bogdanova}} \regalia{\mem{fwinter@yandex.ru}} \position{Postgraduate Student} \dept{Dept. of Higher Mathematics \& Applied computer Science}}





\rukeywords{задача регуляризации, метод Качмажа, регуляризованные нормальные уравнения, уравнения Эйлера}



\ruabstract{Данная статья посвящена модификации итерационного варианта блочного алгоритма Качмажа для решения задачи регуляризации, который является одним из достаточно эффективных методов для задач большой размерности. Важной характеристикой итерационных методов является скорость сходимости, которая зависит от числа обусловленности исходной задачи. Основным недостатком многих итерационных методов является большое число обусловленности, а у методов, основанных на нормальных уравнениях, число обусловленности системы равно квадрату числа обусловленности исходной задачи. В настоящее время для повышения скорости сходимости итерационных методов используются различные типы предобуславливателей, позволяющие снизить число обусловленности системы уравнений. Недостатками данного подхода являются высокая вычислительная сложность, а также отсутствие универсального предобуславливателя, который мог бы применяться для любого итерационного метода. Одним из эффективных подходов для повышения скорости сходимости метода применение использование блочного варианта используемого метода. В связи с этим в данной работе предлагается оригинальная модификация блочного метода Качмажа для задачи регуляризации, которая позволит уменьшить вычислительную сложность и таким образом повысить скорость сходимости алгоритма. В~статье приводится доказательство сходимости предложенного варианта блочного метода Качмажа.

}

%

\entitle{Block regularization Kaczmarz method}

%

\enauthor{E.~Yu.~Bogdanova}{B\,o\,g\,d\,a\,n\,o\,v\,a~E.~Yu.}

%

%

%\enginstitute{\engsamgtu.}% Место работы каждого из авторов на англ. языке

%% Если несколько, то так  \enginstitute{ Организация 1 \and Организация 2  \and Организация 3}

%



%

\enabstract{This article focuses on the modification of the iterative version of Kaczmarz block algorithm for solving the problem of regularization, which is a fairly effective method for large-scale problems. An important characteristic of iterative methods is the speed of convergence, which depends on the condition number of the original problem. The main drawback of many iterative methods is the large condition number, while methods based on normal equations  have the condition number of the system  equal to the square of the condition number of the original problem. At the present time to increase the speed of convergence of iterative methods different types of preconditioners are used reducing the condition number of the system. The disadvantages of this approach is manifested in high computational complexity and the lack of universal preconditioner, which could be applied to any iterative method. One of the most effective approaches for improving the convergence rate of the method is to use a block variant of the method used. In this regard, in this paper we propose a modification of the original block Kaczmarz method for the regularization of the problem, which will reduce the computational complexity, and thus increase the rate of convergence of the algorithm. The article provides a detailed derivation of the proposed modification of the method and the proof of the convergence of the proposed variant of the block Kaczmarz method.

}

%

\enkeywords{regularization problem, Kaczmarz method, regularized normal equations, Euler equations}



\date{17/V/2016} % Дата поступления статьи в редакцию.

\revisiondate{18/VII/2016} % В окончательном виде

\accepteddate{09/IX/2016}



%%%%%%%%%%%%%%%%%%%%%%%%%%%%%%%%%%%%%%%%%%%%%%%%%%%%%%%%%%%%%%%%%%%%%%%%%%%%%%%%%%%%





\makerutitle



\Section[n]{Введение}

Некорректные и плохо обусловленные задачи возникают в различных приложениях при решении интегральных уравнений первого рода и~частных производных, а также в задачах математической физики и~т.~д. 

Для решения некорректных задач существует ряд подходов, одним из которых является метод А.~Н.~Тихонова [\citen{Tihonov}], основанный на решении задачи

\begin{equation}

\min\limits_{u\in \mathbb  R^{n}} \{\|Au-f\|^{2}+\alpha\|u\|^{2}\},

\label{Regularization}

\end{equation}

где $A\in \mathbb R^{m\times n}$, $f\in \mathbb R^{m}$ , $\alpha>0$ "--- параметр регуляризации, 

$\|\cdot\|=\|\cdot\|_{2}$ "--- евклидова норма.



Итерационные алгоритмы для решения задачи \eqref{Regularization} основаны на решении уравнения Эйлера (регуляризованные нормальные уравнения)

\begin{equation}

(A^{\top}A+\alpha E_{n})u=A^{\top}f,

\label{Euler}

\end{equation}

где  $E_{n}$ "--- единичная матрица порядка $n$.



Важной характеристикой итерационных методов является скорость сходимости, которая зависит от числа обусловленности исходной задачи. 

Плохая обусловленность приводит к снижению скорости сходимости итерационных методов. 

Спектральное число обусловленности задачи \eqref{Euler} равно квадрату числа обусловленности исходной задачи \eqref{Regularization}, поэтому у данной задачи \eqref{Euler} медленная скорость сходимости. 

Кроме этого, итерационные методы для решения поставленной задачи расходятся, и поэтому практически невозможно с~помощью них решить задачу (\ref{Regularization}).



Так как итерационные алгоритмы могут иметь крайне медленную сходимость, существуют различные подходы для ее увеличения, одним из которых является предобуславливатель [\citen{Gill}--\citen{Benzi3}], позволяющий снизить число обусловленности системы уравнений. 

Для итерационного метода типа проекционного алгоритма последнее время используется рандомизированый подход [\citen{Bocsan}--\citen{GdanovS}]. 

Однако эти два способа не решают проблему повышения скорости решения задачи. 

Не для всех задач найдены предобуславлеватели. 

При условии, что они найдены, их использование приводит к сильному увеличению вычислительной сложности алгоритма, от чего  снижается эффект от итерационного метода.



Одним из самых наиболее эффективных итерационных методов решения задачи является разработка блочного варианта \cite{Ivanov}.







Для решения системы (\ref{Euler}) рассмотрим блочную форму метода Качмажа [\citen{Vasilchenko}—\citen{Strohmer}], которая основана на проекционном итерационном алгоритме, предложенном в работе \cite{Kaczmarz} польским математиком С.~Качмажем. Применив блочный метод Качмажа к расширенной системе (\ref{Regularization}), предложим оригинальную модификацию данного метода, полученную по аналогии с~работой [\citen{Ivanov}]. 





%Для решения системы (\ref{Euler}) рассмотрим блочную форму метода Качмажа [\citen{Vasilchenko}--\citen{Strohmer}], которая основана на проекционном итерационном алгоритме, предложенном в работе \cite{Kaczmarz} польским математиком С.~Качмажем. 

%Применив блочный метод Качмажа к расширенной системе \eqref{Regularization}, предложим оригинальную модификацию данного метода.



Целью данной работы является получение эффективной вычислительной модификации блочного метода Качмажа для задачи регуляризации, что позволит снизить его вычислительную сложность и~тем самым повысить скорость сходимости алгоритма.



Предполагается, что параметр регуляризации $\alpha$ известен, поэтому в~данной статье не рассматривается его выбор. 

В~настоящее время существует много подходов по его определению, одним из которых является метод, предложенный В.~А.~Морозовым и С.~Ф.~Гилязовым \cite{Morozov}, когда система совместна. 

Также предлагается множество устойчивых правил выбора параметра регуляризации \cite{Hamarik, Dolishniy}, которые являются стабильными с~точки зрения оценки уровня производительности. 

Существуют различные правила отбора параметров регуляризации, не требующие априорной информации \cite{Zhdanov2}.



В~заключение  приведем доказательство сходимости предлагаемого блочного варианта метода Качмажа для задачи регуляризации.



\smallskip





\Section[n]{Постановка задачи}

Регуляризованная нормальная система уравнений \eqref{Euler} может быть представлена в виде матричного уравнения

\begin{equation}

\begin{pmatrix} \omega E_{m} & A \\ A^{\top} & -\omega E_{n} \end{pmatrix}

\begin{pmatrix} y \\ u \end{pmatrix} =

\begin{pmatrix} f \\ 0 \end{pmatrix}

\Longleftrightarrow \widetilde{A_{\omega}}z=\widetilde{f},

\label{Matrix}

\end{equation}

где $\omega=\sqrt{\alpha}$, $E_{m}\in \mathbb R^{m\times m}$, $E_{n}\in \mathbb R^{n\times n}$.



Матрица $\widetilde{A_{\omega}}$ системы \eqref{Matrix} не сингулярная для $\alpha >0$ и~имеет единственное решение "--- вектор $z_{\ast}=\begin{pmatrix} y^{\top}_{\ast} & u^{\top}_{\ast} \end{pmatrix}^{\top}$, 

где $u_{\ast}=(A^{\top}A+\alpha E_{n})^{-1}A^{\top}f$, $y_{\ast}=\omega^{-1}r_{\ast}$, $r_{\ast}=f-Au_{\ast}$.



\smallskip



\Section[n]{Модификация блочного метода Качмажа для задачи регуляризации}

Представим решение системы \eqref{Matrix} в виде системы уравнений

\begin{gather}

(\omega E_{m} | A)z=f,

\label{High}

\\

(A^{\top} | -\omega E_{n})z=0.

\label{Bottom}

\end{gather}



Воспользуемся блочным методом Качмажа:

\begin{gather}

z^{p}_{1}=z^{p}_{0}+(\omega E_{m}|A)^{+}

(f-(\omega E_{m}|A)z^{p}_{0}),

\label{High6}

\\

z^{p}_{t+1}=z^{p}_{t}-B^{p}_{t},

\label{Bottom7}

\end{gather}

где $B^{p}_{t}=(A^{\top}_{p}|-\omega E^{p}_{t})^{+}(A^{\top}_{p}|-\omega E^{p}_{t})z^{p}_{t}$, $p=1,2,\ldots,s$, $z^{p}_{0}=z^{p-1}_{n+1}$, $t=1, 2, \ldots$, $E^{p}_{t}\in \mathbb R^{n_{p}\times n}$ "--- $p$ блок единичной матрицы $E_{n}$, т.е.

$E_{n}=(E^{1}_{n}, E^{2}_{n}, \ldots, E^{s}_{n})$,

$A=(A_{1}, A_{2}, \ldots, A_{s})$,

$A_{p}\in \mathbb  R^{m\times n_{p}}$,

количество блоков зависит от размерности матрицы $A$ и количества строк в блоках,

$\sum _{p=1}^{s}n_{p}=n$, $A^{+}$ "--- псевдообратная матрица к $A$.



Уравнение \eqref{Bottom7} представим в виде рекуррентных уравнений

\begin{gather}

y^{p}_{t+1}=y^{p}_{t}-A_{p}B^{p}_{t},

\label{High8}

\\

u^{p,t}_{t+1}=u^{p,t}_{t}+\omega (E^{p}_{n})^{\top}B^{p}_{t},

\label{High9}

\end{gather}

где $B^{p}_{t}=(G_{p}G^{\top}_{p})^{-1}G_{p}(y^{p}_{t}|u^{p}_{t})^{\top}$, $G_{p}=(A^{\top}_{p}|-\omega E^{p}_{n})$.



Поскольку $y=\omega^{-1}r$, где $r=f-Au$, рекуррентные уравнения \eqref{High8}, \eqref{High9} могут быть преобразованы следующим образом:

\begin{equation}

r^{p}_{t+1}=r^{p}_{t}-A_{p}C^{p}_{t},

\label{High10}

\end{equation}

\begin{equation}

u^{p,t}_{t+1}=u^{p,t}_{t}+(E^{p}_{n})^{\top}C^{p}_{t},

\label{High11}

\end{equation}

где $C^{p}_{t}=(D_{p}D^{\top}_{p})^{-1}D_{p}(y^{p}_{t}|u^{p}_{t})^{\top}$, $D_{p}=(A^{\top}_{p}|-\alpha E^{p}_{n})$.



Покажем, что %не нужно 

можно не

использовать первое рекуррентное уравнение \eqref{High6} в~процессе итерации, если выполнены дополнительные условия "--- совпадают с начальными условиями $u_{0}=u^{1}_{0}$ и $r_{0}=r^{1}_{0}$. 

Уравнение \eqref{High} используется только для согласования c начальными условиями $u_{0}$ и $r_{0}$:

\begin{gather}

r_{k+1}=r_{k}-A_{j(k)}C_{k},

\label{High12}

\\

u_{k+1}=u_{k}+(E^{j(k)}_{n})^{\top}C_{k},

\label{High13}

\end{gather}

где 

$C_{k}=(D_{j(k)}D^{\top}_{j(k)})^{-1}D_{j(k)}(y_{k}|u_{k})^{\top}$, 

$D_{j(k)}=(A^{\top}_{j(k)}|-\alpha E^{j(k)}_{n})$,

$k=0, 1, 2,\ldots$,

$j(k)=k \mathop{\mathrm{mod}}(l)+1,$

$l$ "--- количество блоков, следовательно, $\{j(k)\}^{\infty}_{k=0}$ "--- периодическая последовательность вида 

$\{1, 2, \ldots, l, 1, 2, \ldots\}.$



Введем вектор $\theta _{k}= (r^{\top}_{k},u^{\top}_{k})^{\top}$, с помощью которого 

рекуррентные уравнения  \eqref{High12}, \eqref{High13} можно записать в виде одного рекуррентного уравнения:

\begin{equation}

\theta_{k+1}=\theta_{k}+(A^{\top}_{j(k)}|-E^{j(k)}_{n})^{\top}C_{k},

\label{High14}

\end{equation}

где $C_{k}=(D_{j(k)}D^{\top}_{j(k)})^{-1}D_{j(k)}\theta_{k}$,

$D_{j(k)}=(A^{\top}_{j(k)}|-\alpha E^{j(k)}_{n})$, $k=0, 1, 2, \ldots$, $\theta_{0}$ "--- вектор начальных значений.



\smallskip



{\small \sc Теорема.} \it

Пусть вектор начальных условий $\theta_{0}=(r^{\top}_{0}, u^{\top}_{0})^{\top}$ удовлетворяет условию согласования $r_{0}=f-Au_{0}.$

Тогда последовательность $\theta_{k},$ формируемая рекуррентным уравнением \rm \eqref{High14}, \it сходится к~$\theta_{\ast}{:}$ $\theta_{k} \rightarrow \theta_{\ast}$ при $k \rightarrow \infty,$ $\theta_{\ast}=(r^{\top}_{\ast},u^{\top}_{\ast})^{\top}.$



\rm



\smallskip



{\it Д\,о\,к\,а\,з\,а\,т\,е\,л\,ь\,с\,т\,в\,о.}

Доказательство проведем методом математической индукции. 

Из условия согласования начальных значений $r_{0}=f-Au_{0}$ следует, что условие согласования выполняется при любом $k\geq 0$:

\begin{equation}

r_{k}=f-Au_{k},

\label{High15}

\end{equation}

где $r_{k}$ и $u_{k}$ рассчитываются из рекуррентных уравнений \eqref{High9}, \eqref{High10}.



Для $k=1$ из \eqref{High9}, \eqref{High10} получаем

\begin{multline*}

f-Au_{1}=f-A(u_{0}+(E^{1}_{n})^{\top}C_{0}) = \\ =(f-Au_{0})-A(E^{1}_{n})^{\top}C_{0}=r_{0}-A_{1}C_{0}=r_{1}.

\end{multline*}



Предположим, что \eqref{High15} выполняется для некоторого произвольного $k=\nu$.

Покажем, что \eqref{High15} выполняется для $k=\nu+1$:

\begin{multline*}

f-Au_{\nu+1}=f-A(u_{\nu+1}+(E^{\nu+1}_{n})^{\top}C_{\nu})= \\ =(f-Au_{\nu})-A_{\nu}(E^{\nu+1}_{n})^{\top}C_{\nu}=r_{\nu}-A_{\nu+1}C_{\nu}=r_{\nu+1}.

\end{multline*}

Из этого следует, что \eqref{High15} выполняется для любого $k>0$.



Таким образом, по методу математической индукции для произвольного начального вектора $u_{0}$, $\theta _{k} \rightarrow \theta _{\ast}$, при $k \rightarrow \infty$, $\theta _{\ast}=(r^{\top}_{\ast},u^{\top}_{\ast})$.



В работе \cite{Zhdanov}  показано, что при выполнении теоремы нет необходимости использовать в алгоритме итерационное уравнение~\eqref{High6}. 

Таким образом, блочный алгоритм Качмажа для задачи регуляризации сходится.~$\square$



\smallskip



\Section[n]{Выводы}

В данной статье предложен эффективный итерационный вариант блочного алгоритма Качмажа для решения задачи регуляризации. Данная оригинальная модификация метода позволяет уменьшить вычислительную сложность и повысить скорость сходимости алгоритма для решения некорректных и плохо обусловленных задач, которые возникают в различных приложениях при решении интегральных уравнений первого рода, частных производных, а так же в задачах математической физики.









\thanks{{\bf Декларация о финансовых и других взаимоотношениях.}

Исследование не имело спонсорской поддержки.

Автор несет полную ответственность за предоставление окончательной версии рукописи в печать.

Окончательная версия рукописи была одобрена автором.

Автор не получал гонорар за статью.

}







\thanks{{\bf ORCID}



{\bf  Екатерина Юрьевна Богданова:}    \url{http://orcid.org/0000-0002-8687-3173}



}



\RusRef

\begin{thebibliography}{19}



\RBibitem{Tihonov}

\by Тихонов~А.~Н., Арсений~В.~Я.

\book Методы решения некорректных задач

\publaddr М.

\publ Наука

\yr 1979

\totalpages 284

\zmath{https://zbmath.org/?q=an:0499.65030}



\Bibitem{Gill}

\crossref{10.1137/0613022}

\paper Preconditioners for indefinite systems arising in optimization

\by Gill~P.~E., Murray~W., Saunders~M.~A.

\jour SIAM. J. Matrix Anal. Appl.

\yr 1992

\vol 13

\issue 1

\pages 292--311



\Bibitem{Benzi}

\crossref{10.1006/jcph.2002.7176}

\paper Preconditioning Techniques for Large Linear Systems: A Survey

\by Benzi~M.

\jour J. Comput. Phys.

\yr 2002

\vol 182

\issue 2

\pages 418–477



\Bibitem{Benzi2}

\crossref{10.1016/S0168-9274(98)00118-4}

\paper A comparative study of sparse approximate inverse preconditioners

\by Benzi~M., Tuma~M.

\jour Appl. Numer. Math.

\yr 1999

\vol 30

\issue 1--2

\pages 305--340



\Bibitem{Bergamaschi}

\crossref{10.1002/(SICI)1099-1506(200004/05)7:3<99::AID-NLA188>3.0.CO;2-5}

\paper Aproximate inverse preconditioning in the parallel solution of sparse eigenproblems

\by Bergamaschi~L., Pini~G., Sartoretto~F.

\jour Numer. Linear Algebra Appl.

\yr 2000

\vol 7

\issue 3

\pages 99--116



\Bibitem{Benzi3}

\paper Numerical experiments with parallel orderings for ILU preconditioners

\by Benzi~M., Joubert~W.~D., Mateescu~G.

\jour ETNA. Electronic Transactions on Numerical Analysis

\yr 1999

\vol 8

\pages 88--114

\elink \url{http://eudml.org/doc/119978}



\Bibitem{Bocsan}

\paper Convergence of iterative methods for solving random operator equations

\by Boc\k{s}an~Gh.

\jour J.~Nonlinear Sci. Appl.

\yr 2013

\vol 6

\issue 1

\pages 2--6

\elink \url{http://www.tjnsa.com/includes/files/articles/Vol6_Iss1_2--6_Convergence_of_iterative_methods_fo.pdf}



\Bibitem{Gower}

\paper Randomized Iterative Methods for Linear Systems

\by Gower~R.~M., Richt\'arik~P.

\jour SIAM. J. Matrix Anal. \& Appl.

\yr 2015

\vol 36

\issue 4

\pages 1660–1690

\arxiv \href{https://arxiv.org/abs/1506.03296}{1506.03296} [math.NA]



\RBibitem{GdanovS}

\elib{http://elibrary.ru/item.asp?id=24373929}

\paper Параллельная реализация рандомизированного регуляризованного алгоритма Качмажа

\by Жданов~А.~И., Сидоров~Ю.~В.

\jour Компьютерная оптика

\yr 2015

\vol 39

\issue 4

\pages 536--541

\scopus{2-s2.0-84945194344}

\mathnet{http://mi.mathnet.ru/co15}

\crossref{10.18287/0134-2452-2015-39-4-536-541}



\Bibitem{Ivanov}

\paper Kaczmarz algorithm for Tikhonov regularization problem

\by Ivanov~A.~A., Zhdanov~A.~I.

\jour Applied Mathematics E-Notes 

\yr 2013

\vol 13

\pages 270--276

\elink \url{http://www.math.nthu.edu.tw/~amen/2013/1302252(final).pdf}



\RBibitem{Vasilchenko}

\by Васильченко~Г.~П., Светлаков~А.~А. 

\paper Проекционный алгоритм решения систем линейных алгебраических уравнений большой размерности 

\jour Ж. вычисл. матем. и матем. физ.

\yr 1980

\issue 1

\pages 3--10

\mathnet{http://mi.mathnet.ru/zvmmf5240}

\mathscinet{http://www.ams.org/mathscinet-getitem?mr=564772}

\zmath{https://zbmath.org/?q=an:0446.65012}

%\transl

%\paper A projection algorithm for solving systems of linear algebraic equations of high dimensionality

%\by Vasil'chenko~G.~P., Svetlakov~A.~A.

%\jour U.S.S.R. Comput. Math. Math. Phys.

%\yr 1980

%\issue 1

%\vol 20

%\pages 1--8



\Bibitem{Tanabe}

\paper Projection method for solving a singular system of linear equation and its applications

\by Tanabe~K.

\jour Numer. Math.

\yr 1971

\issue 3

\vol 17

\pages 203--214

\crossref{10.1007/BF01436376}



\Bibitem{Strohmer}

\crossref{doi:10.1007/s00041-008-9030-4}

\paper A randomized Kaczmarz algorithm for linear systems with exponential convergence

\by Strohmer~T.~A., Vershynin~R.

\jour  J. Fourier Anal. Appl.

\yr 2009

\vol 15

\pages 262--278



\Bibitem{Kaczmarz}

\paper Angen\"aherte Aufl\"osung von Systemen linearer Gleichungen

\by Kaczmarz~S.

\jour Bull. Int. Acad. Polon. Sci. A 

\vol 35

\yr 1937

\pages 335--357

\zmath{https://zbmath.org/?q=an:0017.31703|an:63.0524.02}

\elink \url{http://jasonstockmann.com/Jason_Stockmann/Welcome_files/kaczmarz_english_translation_1937.pdf}





\Bibitem{Morozov}

\by Morozov~V.~A.

\book Methods of Solving Incorrectly Posed Problems

\publaddr New York

\publ Springer Verlag

\yr 1984

\crossref{10.1007/978-1-4612-5280-1}.

\totalpages  xviii+257 pp \nofrills

\zmath{https://zbmath.org/?q=an:0658.65002}



\Bibitem{Hamarik}

\paper A family of rules for parameter choice in Tikhonov regularization of ill–posed problems with inexact noise level

\by H\"amarik~U., Palm~R., Raus~T.

\jour Comput. Appl. Math.

\yr 2012

\issue 8

\vol 236

\pages 2146--2157

\crossref{10.1016/j.cam.2011.09.037}



\RBibitem{Dolishniy}

\by Долишний~В.~В., Жданов~А.~И.

\paper Вычисление параметра регуляризации методом перекрестной значимости на основе эквивалентных нормальных расширенных систем

\serial Матем. моделирование и краев. задачи

\yr 2010

\pages 52--55

\inbook Труды седьмой Всероссийской научной конференции с международным участием

\bookinfo Информационные технологии в математическом моделировании. \rm Часть~4

\procinfo 3–6 июня 2010~г.

\yr 2010

\publ СамГТУ

\publaddr Самара

\mathnet{http://mi.mathnet.ru/mmkz1720}





\RBibitem{Zhdanov2}

\by Жданов~А.~И.

\paper Оптимальная регуляризация решений приближенных стохастических систем линейных алгебраических уравнений

\jour Ж. вычисл. матем. и матем. физ.

\yr 1990

\vol 30

\issue 10

\pages 1588--1593

\mathnet{http://mi.mathnet.ru/zvmmf3194}

\mathscinet{http://www.ams.org/mathscinet-getitem?mr=1085712}

\zmath{https://zbmath.org/?q=an:0717.62052}

%\transl

%\jour U.S.S.R. Comput. Math. Math. Phys.

%\yr 1990

%\vol 30

%\issue 5

%\pages 224--229

%\paper Optimal regularization of solutions of approximate stochastic systems of linear algebraic equations

%\by Zhdanov~A.~I.

%\crossref{10.1016/0041-5553(90)90180-Z}











\RBibitem{Zhdanov}

\by Жданов~А.~И.

\paper Метод расширенных регуляризованных нормальных уравнений

\jour Ж. вычисл. матем. и матем. физ.

\yr 2012

\vol 52

\issue 2

\pages 205--208

\mathnet{http://mi.mathnet.ru/zvmmf9648}

\mathscinet{http://www.ams.org/mathscinet-getitem?mr=2953308}

\zmath{https://zbmath.org/?q=an:06057653}

\adsnasa{http://adsabs.harvard.edu/cgi-bin/bib_query?2012CMMPh..52..194Z}

\elib{http://elibrary.ru/item.asp?id=17353053}

%\transl

%\jour Comput. Math. Math. Phys.

%\yr 2012

%\vol 52

%\issue 2

%\pages 194--197

%\paper The method of augmented regularized normal equations

%\by Zhdanov~A.~I.

%\crossref{10.1134/S0965542512020169}

%\isi{http://gateway.isiknowledge.com/gateway/Gateway.cgi?GWVersion=2&SrcApp=PARTNER_APP&SrcAuth=LinksAMR&DestLinkType=FullRecord&DestApp=ALL_WOS&KeyUT=000303535300002}

%\elib{http://elibrary.ru/item.asp?id=17977743}

%\scopus{http://www.scopus.com/record/display.url?origin=inward&eid=2-s2.0-84857590211}









\RBibitem{Gdanov}

\by Жданов~А.~И.

\paper Об одной модификации итерационного алгоритма Качмажа

\serial Матем. моделирование и краев. задачи

\inbook Труды седьмой Всероссийской научной конференции с международным участием

\bookinfo Информационные технологии в математическом моделировании. \rm Часть~4

\procinfo 3–6 июня 2010~г.

\yr 2010

\pages 75--77

\publ СамГТУ

\publaddr Самара

\mathnet{http://mi.mathnet.ru/mmkz1725}



\end{thebibliography}



\RuDate



\bigskip



\makeentitle





\thanks{{\bf Declaration of Financial and Other Relationships.}

The research has not had any sponsorship.

The author is absolutely responsible for submitting the final manuscript in print.

The author has approved the final version of manuscript.

The author has not received any fee for the article.

}





\thanks{{\bf ORCID}



{\bf  Ekaterina Yu. Bogdanova:}    \url{http://orcid.org/0000-0002-8687-3173}



}









\EngRef

\begin{thebibliography}{19}



\Bibitem{Tihonov:en}

\lang In Russian

\by Tikhonov~A.~N., Arsenii~V.~Ya.

\book Metody resheniia nekorrektnykh zadach \rm [Methods for the solution of ill-posed problems]

\publaddr Moscow

\publ Nauka

\yr 1979

\totalpages 284



\Bibitem{Gill:en}

\crossref{10.1137/0613022}

\paper Preconditioners for indefinite systems arising in optimization

\by Gill~P.~E., Murray~W., Saunders~M.~A.

\jour SIAM. J. Matrix Anal. Appl.

\yr 1992

\vol 13

\issue 1

\pages 292--311



\Bibitem{Benzi:en}

\crossref{10.1006/jcph.2002.7176}

\paper Preconditioning Techniques for Large Linear Systems: A Survey

\by Benzi~M.

\jour J. Comput. Phys.

\yr 2002

\vol 182

\issue 2

\pages 418–477



\Bibitem{Benzi2:en}

\crossref{10.1016/S0168-9274(98)00118-4}

\paper A comparative study of sparse approximate inverse preconditioners

\by Benzi~M., Tuma~M.

\jour Appl. Numer. Math.

\yr 1999

\vol 30

\issue 1--2

\pages 305--340



\Bibitem{Bergamaschi:en}

\crossref{10.1002/(SICI)1099-1506(200004/05)7:3<99::AID-NLA188>3.0.CO;2-5}

\paper Aproximate inverse preconditioning in the parallel solution of sparse eigenproblems

\by Bergamaschi~L., Pini~G., Sartoretto~F.

\jour Numer. Linear Algebra Appl.

\yr 2000

\vol 7

\issue 3

\pages 99--116



\Bibitem{Benzi3:en}

\paper Numerical experiments with parallel orderings for ILU preconditioners

\by Benzi~M., Joubert~W.~D., Mateescu~G.

\jour ETNA. Electronic Transactions on Numerical Analysis

\yr 1999

\vol 8

\pages 88--114

\elink \url{http://eudml.org/doc/119978}



\Bibitem{Bocsan:en}

\paper Convergence of iterative methods for solving random operator equations

\by Boc\k{s}an~Gh.

\jour J.~Nonlinear Sci. Appl.

\yr 2013

\vol 6

\issue 1

\pages 2--6

\elink \url{http://www.tjnsa.com/includes/files/articles/Vol6_Iss1_2--6_Convergence_of_iterative_methods_fo.pdf}



\Bibitem{Gower:en}

\paper Randomized Iterative Methods for Linear Systems

\by Gower~R.~M., Richt\'arik~P.

\jour SIAM. J. Matrix Anal. \& Appl.

\yr 2015

\vol 36

\issue 4

\pages 1660–1690

\arxiv \href{https://arxiv.org/abs/1506.03296}{1506.03296} [math.NA]



\Bibitem{GdanovS:en}

\lang In Russian

\by Zhdanov~A.~I., Sidorov~Yu.~V.

\paper Parallel implementation of a randomized regularized Kaczmarz's algorithm

\jour      Computer Optics

\yr 2015

\vol 39

\issue 4

\pages 536--541

\crossref{10.18287/0134-2452-2015-39-4-536-541}



\Bibitem{Ivanov:en}

\paper Kaczmarz algorithm for Tikhonov regularization problem

\by Ivanov~A.~A., Zhdanov~A.~I.

\jour Applied Mathematics E-Notes 

\yr 2013

\vol 13

\pages 270--276

\elink \url{http://www.math.nthu.edu.tw/~amen/2013/1302252(final).pdf}



\Bibitem{Vasilchenko:en}

\mathscinet{http://www.ams.org/mathscinet-getitem?mr=564772}

\zmath{https://zbmath.org/?q=an:0446.65012}

\paper A projection algorithm for solving systems of linear algebraic equations of high dimensionality

\by Vasil'chenko~G.~P., Svetlakov~A.~A.

\jour U.S.S.R. Comput. Math. Math. Phys.

\yr 1980

\issue 1

\vol 20

\pages 1--8



\Bibitem{Tanabe:en}

\paper Projection method for solving a singular system of linear equation and its applications

\by Tanabe~K.

\jour Numer. Math.

\yr 1971

\issue 3

\vol 17

\pages 203--214

\crossref{10.1007/BF01436376}



\Bibitem{Strohmer:en}

\crossref{doi:10.1007/s00041-008-9030-4}

\paper A randomized Kaczmarz algorithm for linear systems with exponential convergence

\by Strohmer~T.~A., Vershynin~R.

\jour  J. Fourier Anal. Appl.

\yr 2009

\vol 15

\pages 262--278



\Bibitem{Kaczmarz:en}

\paper Angen\"aherte Aufl\"osung von Systemen linearer Gleichungen

\by Kaczmarz~S.

\jour Bull. Int. Acad. Polon. Sci. A 

\vol 35

\yr 1937

\pages 335--357

\zmath{https://zbmath.org/?q=an:0017.31703|an:63.0524.02}

\elink \url{http://jasonstockmann.com/Jason_Stockmann/Welcome_files/kaczmarz_english_translation_1937.pdf}





\Bibitem{Morozov:en}

\by Morozov~V.~A.

\book Methods of Solving Incorrectly Posed Problems

\publaddr New York

\publ Springer Verlag

\yr 1984

\crossref{10.1007/978-1-4612-5280-1}.

\totalpages  xviii+257 pp \nofrills

\zmath{https://zbmath.org/?q=an:0658.65002}



\Bibitem{Hamarik:en}

\paper A family of rules for parameter choice in Tikhonov regularization of ill–posed problems with inexact noise level

\by H\"amarik~U., Palm~R., Raus~T.

\jour Comput. Appl. Math.

\yr 2012

\issue 8

\vol 236

\pages 2146--2157

\crossref{10.1016/j.cam.2011.09.037}



\Bibitem{Dolishniy}

\lang In Russian

\by Dolishniy~V.~V., Zhdanov~A.~I.

\paper Calculation of the parameter regularization method based on the importance of cross-equivalent normal extension systems

\serial Matem. Mod. Kraev. Zadachi

\yr 2010

\pages 52--55

\inbook Proceedings of the Seventh  All-Russian Scientific Conference with international participation

\bookinfo Information technologies in mathematical modeling. \rm Part~4

\procinfo 3–6 June 2010

\yr 2010

\publ Samara State Technical Univ.

\publaddr Samara





\Bibitem{Zhdanov2:en}

\mathscinet{http://www.ams.org/mathscinet-getitem?mr=1085712}

\zmath{https://zbmath.org/?q=an:0717.62052}

\jour U.S.S.R. Comput. Math. Math. Phys.

\yr 1990

\vol 30

\issue 5

\pages 224--229

\paper Optimal regularization of solutions of approximate stochastic systems of linear algebraic equations

\by Zhdanov~A.~I.

\crossref{10.1016/0041-5553(90)90180-Z}











\Bibitem{Zhdanov:e}

\mathscinet{http://www.ams.org/mathscinet-getitem?mr=2953308}

\zmath{https://zbmath.org/?q=an:06057653}

\adsnasa{http://adsabs.harvard.edu/cgi-bin/bib_query?2012CMMPh..52..194Z}

\elib{http://elibrary.ru/item.asp?id=17353053}

\jour Comput. Math. Math. Phys.

\yr 2012

\vol 52

\issue 2

\pages 194--197

\paper The method of augmented regularized normal equations

\by Zhdanov~A.~I.

\crossref{10.1134/S0965542512020169}

\isi{http://gateway.isiknowledge.com/gateway/Gateway.cgi?GWVersion=2&SrcApp=PARTNER_APP&SrcAuth=LinksAMR&DestLinkType=FullRecord&DestApp=ALL_WOS&KeyUT=000303535300002}

\elib{http://elibrary.ru/item.asp?id=17977743}

\scopus{http://www.scopus.com/record/display.url?origin=inward&eid=2-s2.0-84857590211}









\Bibitem{Gdanov:en}

\lang In Russian

\by  Zhdanov~A.~I.

\serial Matem. Mod. Kraev. Zadachi

\inbook Proceedings of the Seventh  All-Russian Scientific Conference with international participation

\bookinfo Information technologies in mathematical modeling. \rm Part~4

\procinfo 3–6 June 2010

\yr 2010

\publ Samara State Technical Univ.

\publaddr Samara

\paper A modification of the iterative algorithm Kaczmarz

\pages 75--77



\end{thebibliography}



\EngDate



\Russian

\newpage



\normalsize



%

%\end{document}

